\section{Validation Artifacts \& Procedures}

Every scripted artifact below lives in the repository \citep{runline2026repo},
takes fixed inputs, and writes its results to a JSON file plus a short
summary. Anyone with the repository and its pinned environment can rerun
them. Appendix~\ref{app:repro} lists the exact commands, and
Appendix~\ref{app:checklists} holds the manual checklists.

\paragraph{Who did the work.}
The team members wrote the requirements, chose the benchmark and its
thresholds, reviewed every change through pull requests, and are responsible
for the manual checks. Much of the measurement tooling (the benchmark
harness, the crossing audit, the pool comparison and the acceptance and
deployment scripts) was written with AI coding assistants (Claude Code and
Codex), which ran under our direction and whose output we reviewed before
merging. \todo{confirm this disclosure matches the course's generative-AI
policy and how each member used these tools.} The automated measurements in
Section~\ref{sec:evidence} were run on \todo{date of final runs} on a MacBook
running Python~3.13 with the exact package versions pinned in
\texttt{requirements.txt}. The pins matter because the saved map graphs only
load under the library versions that built them.

\subsection{Automated tests}
The suite has 223 Python tests across 24 files and 35 Node tests that load the
real browser scripts (\texttt{web/app.js} and \texttt{web/run.js}) into a
sandbox. Table~\ref{tab:tests} groups them by what they protect. Tests build
small synthetic graphs where an answer can be worked out by hand, for example
a straight crossing of a major road versus a turn that stays on the same
side, and they also run complete requests on the shipped Charlottesville
graph. A GitHub Actions workflow runs both suites on every pull request; it
was added during this validation work (pull request \#21), after we found
that only Vercel's preview build had been checking pull requests. A pull
request that broke tests could previously show as passing.

\begin{table}[ht]
\centering\small
\caption{Automated test suite by area (Python unless noted).}
\label{tab:tests}
\begin{tabularx}{\linewidth}{@{}Y r Y@{}}
\toprule
Area & Tests & What it checks \\
\midrule
Graph building and metrics (\texttt{osm}, \texttt{pbf}, signal builders) & 63 & Simplification keeps geometry and signal nodes; crossing, signal and repetition counts on hand-built cases \\
Pool storage and serving & 34 & Format round trips, graph-hash checks, neighbour reads, corrupt or stale pools fall back to live planning \\
Metro tiling and tile serving & 52 & Each cell owned by exactly one tile; tile loading, cache bounds, metro precompute \\
Scoring, planning and coverage & 22 & Penalties applied correctly; preferences change rank; coverage limits reported \\
Start-to-finish routes & 20 & Snapping limits, unreachable or identical endpoints, forwarding rules \\
API, regional forwarding and errors & 16 & Validation errors return explained 4xx; regional outages return 503 \\
GPX and export & 5 & Well-formed GPX whose track matches the route \\
Other (geometry, elevation, starts, privacy, rehearsal) & 11 & Includes a check that the public benchmark file holds no coordinates or addresses \\
Browser logic (Node) & 35 & Request flow, error display, route switching, GPS recording on synthetic samples \\
\bottomrule
\end{tabularx}
\end{table}

\subsection{Route benchmark and acceptance script}
The benchmark uses eight public Charlottesville starts: the Rotunda, the
centre of the map area, and six points 2.5\,km from the centre at 60 degree
intervals. Each start is requested at 3, 5, 8 and 12 miles, giving 32
requests that return up to 96 routes. Private starts that team members use
for their own runs are kept in a git-ignored file and are excluded from every
reported result. The acceptance script (\texttt{scripts/acceptance.py})
plans each request under the default settings (mixed surface, balanced
elevation), then again as road, trail, flat and hilly, and computes each
criterion exactly as Table~\ref{tab:criteria} states it:

\begin{itemize}
\item \textbf{Distance (US1).} Absolute error of the top-ranked route as a
  percentage of the target.
\item \textbf{Alternatives (US2).} For each pair of returned routes, shared
  edge length divided by the shorter route's length; a request passes when
  some pair overlaps by 80\% or less.
\item \textbf{Interruptions (US3).} Mapped signal events plus major-road
  crossings, compared with a \emph{distance-only baseline} planned on the
  same requests. The baseline is a genuine alternative planner, not the
  production one with a setting changed: its path costs are plain edge
  lengths with no signal, crossing or surface penalties, and it ranks only
  by distance error. Both use Dijkstra's algorithm
  \citep{dijkstra1959note}.
\item \textbf{Preferences (US5, US6).} Paired comparisons on the same start
  and distance. A road or trail request ``moves toward'' its surface when its
  trail share is no higher (road) or no lower (trail) than the mixed
  request's. A pair counts for US5 only if the mixed route contains both
  surface types. For US6 the hilly route must have strictly more elevation
  gain than the flat one. Elevation comes from the Copernicus GLO-90 terrain
  model \citep{copernicus2022dem}, baked into the graph when it is built.
\item \textbf{GPX (US7).} Ten routes are sampled across the four distances.
  Each file is parsed and validated against the official GPX~1.1 schema
  \citep{topografix2004gpx}, and its track must be closed and match the
  route's reported length within 1\%.
\item \textbf{Speed, local (US10).} Wall time per request after the map is
  loaded, for 3 and 5 miles.
\end{itemize}

The script runs twice: once against the live planner alone, and once through
the deployed serving path, which answers from precomputed routes first and
plans live only when none fit. The second run measures what a runner actually
receives. The graph's hash is recorded with the results so they can be tied
to the exact map they came from.

\subsection{Shape and crossing audits}
\texttt{scripts/shape\_audit.py} plans 96 routes from a ring of starts and
records properties that the acceptance criteria do not mention but that make
a route bad to run: roundness, measured as $4\pi A/L^2$ (the Polsby--Popper
score \citep{polsby1991third}, where 1 is a circle and below 0.2 is close to
an out-and-back), whether the route enters a dead end, short revisits, and the
share of its length that repeats.

\texttt{scripts/audit\_crossings.py} provides independent labels for US3. For
each stretch of route that touches a major road, it draws a 15\,m corridor
with at least 50\,m of travel on either side, cuts it along the road's
centreline, and counts a crossing when the route passes from one piece to the
other. Grade-separated passes are excluded. The planner's own count is then
scored against these labels for precision and recall. The labeller never calls
the planner's crossing rule, so the two can disagree, and when they did we
plotted each disagreement and settled it by inspecting the map.

\subsection{Precomputed route (pool) check}
\texttt{scripts/compare\_pools.py} draws 40 starts at random (seed 7) within
6\,km of the Charlottesville centre and requests 3, 5, 6.5, 8 and 12 miles on
all three surfaces: 600 requests. It records whether the precomputed routes
could answer, how long that took, and the same quality metrics as the
benchmark, so answers from pools can be compared with live planning on the
same requests.

\subsection{Deployed-system validation}
\texttt{scripts/validate\_deployed.py} sends requests to the public URL one at
a time with a pause between them, using public landmark starts only, and
caps the total at about 250 requests to stay polite to the shared servers.
For US10 it times 3 and 5 mile requests at ten Charlottesville starts, twice
each, reporting the first request at each start separately as a cold start.
For US11 it requests 3, 5 and 8 miles from public landmarks in each metro
area and checks that each loop is closed, in the distance band, and has a GPX
file whose track matches the route. For US9 it runs ten nearby-start cases
across 1, 3 and 5 mile radii and checks every suggested start against the
radius. For US4 a Node harness feeds real responses captured from production
into the interface's own rendering code and compares every number shown on a
route card with the API's metrics. Separately,
\texttt{scripts/presentation\_check.py} replays ten fixed scenarios, including
error cases and address autocomplete, and confirms each GPX file matches its
drawn route.

\subsection{Manual checks}
Three procedures need a person. Each team member runs the browser checklist
in Appendix~\ref{app:checklists} (US8) in a desktop browser and on a phone,
recording pass, fail or notes for each step. The team opens at least one GPX
file per distance in an outside application and records whether the track
appears correctly (US7). Each member also independently reviews the same
twelve benchmark routes against a five-point checklist covering continuity,
backtracking, unreasonable crossings, map rendering, and agreement between
the drawn route and its metrics. Disagreements are discussed and the outcome
recorded.

\subsection{Consent, privacy and accessibility}
No outside participants took part, so no consent was needed. Privacy shaped
the procedures in three ways. Only public landmarks appear in tracked files
and in this report, and a test fails the suite if coordinates or addresses
appear in the public benchmark file. The application keeps no account,
history or saved requests. A recorded run stays in the page's memory and is
discarded when the page closes \citep{w3c2026geolocation}. Map data is
OpenStreetMap's, used under the Open Database License with attribution
\citep{osm2026copyright}. We did not evaluate accessibility, which is listed
as a gap in Section~\ref{sec:analysis}.
